\section{Introduction}
\label{sec:introduction}

\subsection{The Distributional Problem in Standard ReCom}

Markov chain redistricting has two distinct goals that are often conflated.
The first is \emph{optimisation}: finding the best plan by some criterion
(edge cut, compactness, partisan fairness).
The second is \emph{sampling}: drawing plans from a target distribution so
that ensemble statistics reflect genuine distributional properties of the
plan space.

Standard \recom{}~\citep{deford2021} is widely used for both goals.
At each step, it selects an adjacent pair of districts, merges them, samples
a random spanning tree of the merged region via Wilson's algorithm, and cuts
the tree at a balanced edge.
Every valid balanced cut is accepted; the chain is never rejected.

This 100\% acceptance rate is computationally appealing, but it creates
a distributional bias.
Plans reachable by many spanning tree cuts---those whose district boundaries
happen to lie in highly-connected regions of the census-tract graph---are
visited more often than plans reachable by few cuts.
The stationary distribution of standard \recom{} is therefore \emph{not} the
uniform distribution over valid plans; it is a distribution weighted by the
number of spanning tree cuts that produce each plan.
DeFord, Duchin, and Solomon~\citep{deford2021} acknowledge this but note that
characterising the bias analytically is extremely difficult for real census
graphs.

For optimisation, this bias is harmless or even beneficial---overweighting
geometrically accessible plans can guide the chain toward structurally better
solutions.
But for ensemble analysis---where the goal is to characterise the full space
of valid plans and answer questions like ``what fraction of valid plans
produce 7 Democratic seats in NC?''---the bias is problematic.
Ensemble statistics derived from a biased chain may not represent the space
of all valid plans.

\subsection{Forest ReCom as the Fix}

McCartan and Imai~\citep{mccartan2020} introduce \fr{} to correct this bias.
The key idea is to apply a Metropolis-Hastings correction to the \recom{}
proposal: after generating the proposal from the forward spanning tree
$T_{\mathrm{fwd}}$, draw an independent reverse spanning tree $T_{\mathrm{rev}}$
from the \emph{proposed} district pair, count its valid balanced cuts, and
accept the proposal with probability
\[
  \alpha = \min\!\left(1,\; \frac{c_{\mathrm{fwd}}}{c_{\mathrm{rev}}}\right),
\]
where $c_{\mathrm{fwd}}$ is the number of valid cuts in $T_{\mathrm{fwd}}$
and $c_{\mathrm{rev}}$ is the number of valid cuts in $T_{\mathrm{rev}}$.
If $c_{\mathrm{rev}} = 0$, the proposal is rejected with probability 1.

The intuition is direct: if the forward tree has many more balanced cuts than
the reverse tree, the proposed plan is ``more accessible'' from the current
plan than vice versa; the Metropolis-Hastings ratio downweights this
asymmetry.
Over many steps, the chain converges toward the uniform distribution over valid
plans.

\fr{} does not achieve \emph{exact} detailed balance in the per-step sense
because $T_{\mathrm{rev}}$ is a single spanning tree sample, not the
expectation over all spanning trees.
The ratio $c_{\mathrm{fwd}} / c_{\mathrm{rev}}$ is an unbiased estimator of
the true transition probability ratio $Q(\pi \to \pi') / Q(\pi' \to \pi)$
only in expectation.
We characterise this precisely in Section~\ref{sec:algorithm} and follow
McCartan and Imai~\citep{mccartan2020} §3.2 in calling the resulting guarantee
\emph{expected detailed balance}.

\subsection{Three Contributions}

\begin{enumerate}
  \item \textbf{Implementation.}
        We provide the complete specification and implementation of \fr{} as
        \texttt{ForestRecomChain} in \texttt{redist-ensemble}, including the
        deterministic forward/reverse seed schedule, pair selection with
        reselection, and the acceptance criterion
        (Section~\ref{sec:algorithm}).

  \item \textbf{Comparison.}
        We compare \fr{} against standard \recom{} on NC ($k = 14$), WI
        ($k = 8$), and TX ($k = 38$) across acceptance rate, edge-cut
        distribution, partisan outcomes, and runtime
        (Section~\ref{sec:comparison}).

  \item \textbf{Distributional correctness analysis.}
        We verify convergence to the uniform plan distribution on a small
        synthetic graph where the exact uniform distribution is computable,
        and use a chi-squared test to quantify the improvement over standard
        \recom{} (Section~\ref{sec:distributional}).
\end{enumerate}

\subsection{Paper Structure}

Section~\ref{sec:background} reviews standard \recom{}, the spanning-tree
selection bias, and Wilson's \ust{} algorithm.
Section~\ref{sec:algorithm} gives the formal \fr{} algorithm and proves
expected detailed balance.
Section~\ref{sec:comparison} presents empirical comparisons on real states.
Section~\ref{sec:distributional} establishes distributional correctness on
a synthetic benchmark.
Section~\ref{sec:conclusion} concludes with usage guidance and open
questions.
